% \subsection*{Задание 1: Любовный треугольник}
% \begin{figure}[H]
% 	\centering
% 	\includegraphics[width=\textwidth, height=0.9\textheight, keepaspectratio]{1.pdf}
% 	\caption{Схема алгоритма задачи}
% \end{figure}
% \newpage
% \lstinputlisting[caption={Исходный код программы (C)}]{1.c}
% \noindent \indent
% Данное задание можно решить просто пройдясь по всем элементам массива и сравнить текущий индекс элемента
% с индексом котрый возвращает цепочка из трёх элементов.
% \newpage
% \subsection*{Задание 2: Грегор и игра в пешки}
% \begin{figure}[H]
% 	\centering
% 	\includegraphics[width=\textwidth, height=0.9\textheight, keepaspectratio]{2.pdf}
% 	\caption{Схема алгоритма задачи}
% \end{figure}
% \newpage
% \lstinputlisting[caption={Исходный код программы (C)}]{2.c}
% \noindent \indent
% Это задание можно решить сразу при вводе, проверя прямой ход пешки или наискосок если там стоит вражеская пешка.
% \newpage
% \subsection*{Задание 3: Великий побег Авраама}
% \begin{figure}[H]
% 	\centering
% 	\includegraphics[width=\textwidth, height=0.9\textheight, keepaspectratio]{3.pdf}
% 	\caption{Схема алгоритма задачи}
% \end{figure}
% \newpage
% \lstinputlisting[caption={Исходный код программы (C)}]{3.c}
% \noindent \indent
% Для решения этого задания надо вывести лабиринт в котором k точек старта приводят к выходу, а остальные нет.
% Для этого можно завести счётчик и сводить к выходу влево-вверх пока счётчик меньше k, а иначе сводидть вправо-вниз
% и развернуть последнюю влево.
% \newpage
% \subsection*{Задание 4: Топология сети}
% \begin{figure}[H]
% 	\centering
% 	\includegraphics[width=\textwidth, height=0.9\textheight, keepaspectratio]{4.pdf}
% 	\caption{Схема алгоритма задачи}
% \end{figure}
% \newpage
% \lstinputlisting[caption={Исходный код программы (C)}]{4.c}
% \newpage
% \noindent \indent
% Это задание легко решается, путём проверки графа на степени вершин, то есть для шины - 
% две вершини имеют степень 1, а остальные степень 2. Для кольца - все вершины имеют степень 2.
% Так как уже известно что граф связан, то нет необходимости проверять шину и кольцо. Для звезды -
% одна вершина имеет $m$ степень, а остальные степень 1.
% % \newpage